% Prototype only. Live book macro: metadata/authorship-bars.tex
\documentclass[11pt]{article}
\usepackage[a4paper,margin=28mm]{geometry}
\usepackage{amsmath}
\usepackage[most]{tcolorbox}
\usepackage{xcolor}

\pagestyle{empty}

% TikZ/tcolorbox-based paragraph bars. The right edge can use any TikZ style.
\newtcolorbox{rightbar}[1][]{
  enhanced,
  breakable,
  frame hidden,
  colback=white,
  boxrule=0pt,
  sharp corners,
  left=0pt,
  right=7mm,
  top=0pt,
  bottom=0pt,
  before skip=\baselineskip,
  after skip=\baselineskip,
  borderline east={0.9pt}{0pt}{black,#1}
}

\begin{document}

\section*{TikZ right-bar demo}

Normal paragraph without a bar. It establishes the ordinary text width and spacing used by the examples below.

\begin{rightbar}
This paragraph has a solid bar. It contains ordinary inline mathematics such as
$E=mc^2$ and a tall inline expression
$\displaystyle \sum_{k=1}^{n}\frac{a_k}{1+a_k^2}$.
The bar height is determined from the actual box contents rather than estimated from the number of lines.
\end{rightbar}

\begin{rightbar}[dashed]
This paragraph has a dashed bar. Display mathematics also works:
\[
  F(x)=\int_0^\infty e^{-xt}f(t)\,dt,
  \qquad
  \frac{d}{dx}\log Z(x)=\frac{Z'(x)}{Z(x)}.
\]
Text may continue after the display, and the bar spans the complete region.
\end{rightbar}

\begin{rightbar}[densely dotted,line width=1.2pt]
This one is densely dotted and slightly thicker. The optional argument is passed directly to TikZ, so styles such as \texttt{dashed}, \texttt{densely dotted}, colors, and custom dash patterns can be selected per paragraph.
\end{rightbar}

Normal text continues afterward.

\end{document}
